\chapter*{Abstract}
\addcontentsline{toc}{chapter}{Abstract}

Despite significant advances in artificial intelligence and robotics, and recent progress toward more flexible and open-world-capable systems, many approaches still exhibit brittle behavior when confronted with changes beyond their training distributions, reflecting a continued reliance on closed-world assumptions. In real-world environments, new objects, altered dynamics, and unforeseen task variations arise naturally, and even small changes can render existing knowledge inadequate, exposing a limitation in how current systems represent and reason about their interaction with the world.

In this dissertation, we advance the study of open-world agent--environment interaction by developing the formal, algorithmic, and representational foundations required for agents to operate under novelty. We formalize novelty as a mismatch between an agent’s internal models and the environment and introduce environment models, transformations, and evaluation methodologies for its systematic study.
Building on this foundation, we show how abstractions shape learning and behavior in open-world settings. Symbolic abstractions over actions and tasks expose the structure of execution, enabling failures to be localized and converted into targeted learning problems for rapid adaptation.
Object-centric force-space abstractions represent interaction in terms of object dynamics rather than robot-specific control, enabling generalization in robotic manipulation across objects, embodiments, and physical conditions.

Together, these contributions provide a formal framework for studying and addressing novelty in open-world agent--environment interaction. They show how abstractions enable adaptation and generalization under changing conditions, advancing the development of AI and robotic systems that operate robustly in real-world environments.
